\documentclass[11pt]{article}

\usepackage[letterpaper,top=0.92in,bottom=0.88in,left=1in,right=1in]{geometry}
\usepackage{amsfonts}
\usepackage{amsmath,amssymb,amsthm,mathtools}

\theoremstyle{definition}
\newtheorem{theorem}{Theorem}
\newtheorem{lemma}{Lemma}

\renewcommand{\qedsymbol}{\ensuremath{\blacksquare}}
\setlength{\parindent}{0pt}
\setlength{\parskip}{0.6em}
\allowdisplaybreaks

\title{Adversarial delays with secondary-market prices}
\author{Vinayak Pathak}
\date{August 2026}

\begin{document}

\maketitle

\section{Model and path parameters}
\label{sec:model-and-path-parameters}

The player's action space is $[0,1]$. Fix a horizon $T\ge 2$. The
game is parameterized by two unknown deterministic sequences chosen by
an oblivious adversary,
\[
    \boldsymbol{p}=(p_1,p_2,\ldots),
    \qquad
    \boldsymbol{m}=(m_1,m_2,\ldots),
\]
where $p_t,m_t\in[0,1]$. Let
$\boldsymbol{p}_T:=(p_1,\ldots,p_T)$ and
$\boldsymbol{m}_T:=(m_1,\ldots,m_T)$ denote their horizon-$T$ prefixes.
All horizon-$T$ rewards and path parameters below depend only on these
prefixes, so any finite construction may be extended arbitrarily after
time $T$.
The primary-price sequence defines the delay functions
\[
    d_t(a):=\min\bigl\{\tau\in\mathbb{N}_0:p_{t+\tau}>a\bigr\},
\]
where $\mathbb{N}_0=\{0,1,\ldots\}$, with $d_t(a)=\infty$ if the set is
empty. We use the usual convention $t+\infty=\infty$. For notational
convenience, set $m_\infty:=0$; this value is always multiplied by a zero
execution indicator and never affects a reward.

At time $t$:
\begin{enumerate}
    \item The player posts an offer $A_t\in[0,1]$, possibly using
    randomization and the observations from times $1,\ldots,t-1$.
    \item The environment reveals $p_t\in[0,1]$ and the secondary-market
    price $m_t\in[0,1]$.
    \item For each $s\le t$: if the offer $A_s$ is still outstanding and
    $p_t>A_s$, then $A_s$ is executed. The market maker sells at $A_s$,
    immediately buys at $m_t$, and receives reward $A_s-m_t$.
\end{enumerate}

The reward associated with an offer $a$ submitted at time $t$ is
\[
    r_t(a):=
    \bigl(a-m_{t+d_t(a)}\bigr)
    \mathbb{I}\bigl\{t+d_t(a)\le T\bigr\}.
\]

The benchmark obtained by posting the same offer on every round is
\[
    B_T(\boldsymbol{p}_T,\boldsymbol{m}_T)
    :=
    \sup_{a\in[0,1]}\sum_{t=1}^T r_t(a).
\]

The regret of a possibly randomized policy $\pi$ is
\[
    R_T^\pi(\boldsymbol{p}_T,\boldsymbol{m}_T)
    :=
    B_T(\boldsymbol{p}_T,\boldsymbol{m}_T)
    -
    \mathbb{E}_\pi\left[\sum_{t=1}^T r_t(A_t)\right],
\]
where $\mathbb{I}$ denotes the indicator function and
$\mathbb{E}_\pi$ is only over the player's internal randomization. We
write $R_T$ when the policy and paths are clear. Notice that
$B_T(\boldsymbol{p}_T,\boldsymbol{m}_T)\ge0$, because the offer $a=1$
never executes.

For each $t\le T$, define the suffix maximum
\[
    M_t:=\max_{t\le s\le T}p_s.
\]
Then $M_1\ge M_2\ge\cdots\ge M_T$. By definition,
\[
    t+d_t(a)\le T
\]
if and only if $a<M_t$, and
\[
    \bigl\{a\in[0,1]:1\le d_t(a)\le T-t\bigr\}=[p_t,M_t).
\]

Let $U$ be uniformly distributed on $[0,1]$. Define the primary-price
path complexity
\[
    \mathcal{V}_T
    :=
    \sum_{t=1}^{T-1}
    \mathbb{P}\bigl(1\le d_t(U)\le T-t\bigr)
    =
    \sum_{t=1}^{T-1}(M_t-p_t).
\]

The new issue is that the liquidation price attached to a quote depends
on the date at which that quote first executes.

For $s=1,\ldots,T-1$, define the future secondary-price exposure
\[
    \omega_s(\boldsymbol{m}_T)
    :=
    \max_{s<u\le T}|m_u-m_s|,
\]
and define
\[
    W_T(\boldsymbol{m}_T)
    :=
    \sum_{s=1}^{T-1}\omega_s(\boldsymbol{m}_T).
\]

Throughout, when the underlying paths are clear, we suppress the path
arguments of path-dependent quantities.

The quantity $W_T$ depends only on the realized secondary-price path
$\boldsymbol{m}_T$. It does not depend on the primary-price path, the
learner, or an auxiliary predictor. When the offer at time $s+1$ is
selected, the most recently observed secondary-market price is $m_s$.
Any eventual execution date $u$ of that offer satisfies $u>s$, so
$\omega_s$ bounds the largest secondary-price movement to which the
offer could be exposed. The maximum ranges over all future dates,
whether or not they can be the first execution date of an offer.

\section{Upper bound}
\label{sec:upper-bound}

For every $t=1,\ldots,T$, define
\[
    H_t:=\sum_{s=1}^t\frac1s.
\]

Write $(x)_+:=\max\{x,0\}$.

Set
\[
    c_1:=1,
    \qquad
    c_t:=m_{t-1}\quad(t\ge2),
\]
and define the surrogate reward
\[
    \widetilde r_t(a):=(a-c_t)\mathbb{I}\{a<M_t\}.
\]

Let
\[
    \widetilde B_T
    :=
    \sup_{a\in[0,1]}
    \sum_{t=1}^T\widetilde r_t(a).
\]

For every $t=1,\ldots,T$, define
\[
    C_t:=\sum_{s=1}^t c_s.
\]

\begin{lemma}
\label{lem:surrogate-benchmark}
For every pair of paths
$\boldsymbol{p}_T,\boldsymbol{m}_T$,
\[
    \widetilde B_T
    \le
    \max\left\{
        0,\,
        \max_{1\le k\le T}(kM_k-C_k)
    \right\}.
\]
\end{lemma}

\begin{proof}
For $a\in[0,1]$, let
\[
    N(a):=\max\{1\le t\le T:a<M_t\},
\]
with $N(a)=0$ if the set is empty. Since
$M_1\ge M_2\ge\cdots\ge M_T$, if $N(a)\ge1$, then
\[
    \sum_{t=1}^T\widetilde r_t(a)
    =
    \sum_{t=1}^{N(a)}(a-c_t)
    =
    aN(a)-C_{N(a)}.
\]
Moreover, $a<M_{N(a)}$, and hence
\[
    aN(a)-C_{N(a)}
    <
    N(a)M_{N(a)}-C_{N(a)}
    \le
    \max_{1\le k\le T}(kM_k-C_k).
\]
If $N(a)=0$, the surrogate payoff is zero.
Taking the supremum over $a$ concludes the proof.
\end{proof}

Define the clairvoyant actions
\[
    A_1^\star:=c_1
\]
and, for $t\ge2$,
\[
    A_t^\star
    :=
    c_t+\frac{t-1}{t}(M_{t-1}-c_t)_+.
\]

Let
\[
    S_t^\star:=\sum_{s=1}^t\widetilde r_s(A_s^\star).
\]

\begin{lemma}[Clairvoyant potential bound]
\label{lem:clairvoyant-potential}
For every $t=1,\ldots,T$,
\[
    S_t^\star\ge tM_t-C_t-H_t+1.
\]
Consequently,
\[
    S_T^\star\ge\widetilde B_T-(H_T-1).
\]
\end{lemma}

\begin{proof}
We prove the statement by induction on $t$. For $t=1$,
\[
    S_1^\star
    =
    \widetilde r_1(1)
    =
    0
    \ge
    M_1-C_1-H_1+1
    =
    M_1-1.
\]

Now fix $t\ge2$ and assume
\[
    S_{t-1}^\star
    \ge
    (t-1)M_{t-1}-C_{t-1}-H_{t-1}+1.
\]

There are two cases.
\begin{enumerate}
    \item If $A_t^\star<M_t$, then $M_{t-1}>c_t$, the offer executes by
    time $T$, and
    \[
        \widetilde r_t(A_t^\star)
        =
        \frac{t-1}{t}(M_{t-1}-c_t).
    \]
    Therefore
    \begin{align*}
        S_t^\star
        &\ge
        (t-1)M_{t-1}-C_{t-1}-H_{t-1}+1
        +\frac{t-1}{t}(M_{t-1}-c_t)
        \\
        &=
        \left(t-\frac1t\right)M_{t-1}
        -C_t-H_{t-1}+1+\frac{c_t}{t}
        \\
        &\ge
        tM_t-C_t-H_t+1.
    \end{align*}
    For the last step, we used $M_{t-1}\ge M_t$ and
    \[
        \left(t-\frac1t\right)M_t
        +\frac{c_t}{t}
        +\frac1t
        -tM_t
        =
        \frac{1+c_t-M_t}{t}
        \ge0.
    \]

    \item If $A_t^\star\ge M_t$, then the offer does not execute by time
    $T$ and $\widetilde r_t(A_t^\star)=0$. If
    $M_{t-1}\le c_t$, then
    \[
        tM_t
        \le
        tM_{t-1}
        \le
        (t-1)M_{t-1}+c_t.
    \]
    If $M_{t-1}>c_t$, the case condition gives the same inequality.
    Hence
    \[
        S_t^\star
        =
        S_{t-1}^\star
        \ge
        (t-1)M_{t-1}-C_{t-1}-H_{t-1}+1
        \ge
        tM_t-C_t-H_t+1.
    \]
\end{enumerate}

This completes the induction. To prove the last statement, observe that
every clairvoyant surrogate reward is nonnegative. If the outer maximum
in Lemma~\ref{lem:surrogate-benchmark} is positive, choose $k$ achieving
the inner maximum. Then
\[
    S_T^\star
    \ge
    S_k^\star
    \ge
    kM_k-C_k-H_k+1
    \ge
    \widetilde B_T-(H_T-1).
\]
If the outer maximum is zero,
Lemma~\ref{lem:surrogate-benchmark} gives $\widetilde B_T=0$, while
$S_T^\star\ge0$. The same conclusion follows, completing the proof.
\end{proof}

\begin{lemma}
\label{lem:surrogate-one-sided-lipschitz}
For every $t$ and every $c_t\le x\le y\le1$,
\[
    \widetilde r_t(x)
    \ge
    \widetilde r_t(y)-(y-x).
\]
\end{lemma}

\begin{proof}
If $y<M_t$, then also $x<M_t$, and
\[
    \widetilde r_t(x)
    =
    x-c_t
    =
    \widetilde r_t(y)-(y-x).
\]
If $y\ge M_t$, then $\widetilde r_t(y)=0$, and so
$\widetilde r_t(y)-(y-x)\le0$, whereas
$\widetilde r_t(x)\ge0$, concluding the proof.
\end{proof}

\begin{theorem}[Pathwise upper bound]
\label{thm:pathwise-upper-bound}
The deterministic policy
\[
    A_1=1,
    \qquad
    A_t
    =
    m_{t-1}
    +
    \frac{t-1}{t}(p_{t-1}-m_{t-1})_+
    \quad(t\ge2)
\]
satisfies
\[
    R_T\le H_T+\mathcal{V}_T+2W_T.
\]
\end{theorem}

\begin{proof}
For $t\ge2$, one has $p_{t-1}\le M_{t-1}$, and therefore
\[
    c_t\le A_t\le A_t^\star.
\]
Applying Lemma~\ref{lem:surrogate-one-sided-lipschitz} gives
\[
    \widetilde r_t(A_t)
    \ge
    \widetilde r_t(A_t^\star)-(A_t^\star-A_t).
\]
Moreover,
\begin{align*}
    A_t^\star-A_t
    &=
    \frac{t-1}{t}
    \left[
        (M_{t-1}-c_t)_+
        -
        (p_{t-1}-c_t)_+
    \right]
    \\
    &\le
    \frac{t-1}{t}(M_{t-1}-p_{t-1}),
\end{align*}
because $x\mapsto(x-c_t)_+$ is $1$-Lipschitz. The same
surrogate-reward inequality is trivial at $t=1$, because
$A_1=A_1^\star=1$.

We next compare surrogate and actual rewards. For $t\ge2$,
\[
    \left|r_t(a)-\widetilde r_t(a)\right|
    =
    \left|m_{t-1}-m_{t+d_t(a)}\right|
    \mathbb{I}\{t+d_t(a)\le T\}
    \le
    \omega_{t-1},
\]
because, whenever the offer executes by time $T$, its execution date
$t+d_t(a)$ is strictly larger than $t-1$. At $t=1$,
\[
    r_1(a)-\widetilde r_1(a)
    =
    \bigl(1-m_{1+d_1(a)}\bigr)
    \mathbb{I}\{1+d_1(a)\le T\}
    \le1.
\]
Thus, for every fixed offer $a$,
\[
    \sum_{t=1}^T r_t(a)
    \le
    \sum_{t=1}^T\widetilde r_t(a)+1+W_T,
\]
and therefore
\[
    B_T\le\widetilde B_T+1+W_T.
\]

For the learner, $A_1=1$ never executes and
$r_1(A_1)=\widetilde r_1(A_1)=0$. Applying the preceding bound to the
remaining offers and summing over $t$, we obtain
\begin{align*}
    \sum_{t=1}^T r_t(A_t)
    &\ge
    \sum_{t=1}^T\widetilde r_t(A_t)-W_T
    \\
    &\ge
    S_T^\star
    -
    \sum_{t=2}^T
    \frac{t-1}{t}(M_{t-1}-p_{t-1})
    -
    W_T
    \\
    &\ge
    \widetilde B_T-(H_T-1)
    -
    \sum_{s=1}^{T-1}
    \frac{s}{s+1}(M_s-p_s)
    -
    W_T
    \\
    &\ge
    B_T-H_T
    -
    \sum_{s=1}^{T-1}
    \frac{s}{s+1}(M_s-p_s)
    -
    2W_T
    \\
    &\ge
    B_T-H_T-\mathcal{V}_T-2W_T.
\end{align*}

The third inequality uses
Lemma~\ref{lem:clairvoyant-potential}, the fourth uses the preceding
benchmark comparison, and the last one uses the definition of
$\mathcal{V}_T$. Rearranging concludes the proof.
\end{proof}

\section{A matching lower bound for the new term}
\label{sec:a-matching-lower-bound-for-the-new-term}

Although $W_T$ maximizes over all future dates, the next construction
shows that its linear contribution is unavoidable when one terminal
secondary-price movement can affect every earlier offer.

\begin{theorem}[Necessity of $W_T$]
\label{thm:necessity-w}
Fix $T \ge 2$. For every possibly randomized learner and every
\[
    0 \le w \le \frac{T-1}{4},
\]
there exists a deterministic oblivious pair of sequences
$(\boldsymbol{p},\boldsymbol{m})$ such that
\[
    W_T=w,
    \qquad
    \mathcal{V}_T=\frac{w}{4T},
\]
and
\[
    R_T\ge\frac{3}{8}w.
\]
Thus no algorithm can guarantee $o(W_T)$ regret uniformly, even on
sequences satisfying
\[
    \mathcal{V}_T=O\left(\frac{W_T}{T}\right).
\]
\end{theorem}

\begin{proof}
Set
\[
    q:=\frac{3}{4},
    \qquad
    c:=\frac{1}{2},
    \qquad
    \delta:=\frac{w}{T-1},
    \qquad
    \epsilon:=\frac{\delta}{4T}.
\]
Use the primary-price sequence
\[
    p_t=q \quad (t<T),
    \qquad
    p_T=q+\epsilon.
\]
Consider two possible secondary-price sequences. In both worlds,
\[
    m_t=c \quad (t<T),
\]
whereas
\[
    m_T=
    \begin{cases}
        c-\delta & \text{in world L}\\
        c+\delta & \text{in world H}
    \end{cases}.
\]
Put the uniform prior on the two worlds. Before choosing $A_T$, the
learner has observed the same history in both worlds. Hence every
learner action is independent of the world.

For any round, the learner's prior-expected reward is at most
\[
    q+\epsilon-c=\frac{1}{4}+\epsilon.
\]
Indeed, if $a<q$, it executes immediately and receives
$a-c\le q-c$. If $q\le a<q+\epsilon$, it executes at $T$ and its
prior-expected reward is $a-c\le q+\epsilon-c$. If
$a\ge q+\epsilon$, it never executes. Thus the learner's expected
cumulative reward is at most
\[
    T\left(\frac{1}{4}+\epsilon\right).
\]

In world L, the fixed comparator may quote just below $q+\epsilon$,
wait until $T$, and obtain a payoff arbitrarily close to
\[
    T\left(\frac{1}{4}+\epsilon+\delta\right).
\]
In world H, the comparator may quote just below $q$. It executes
immediately on the first $T-1$ rounds and at $T$ on the final round,
obtaining a payoff arbitrarily close to
\[
    (T-1)\frac{1}{4}
    +\left(\frac{1}{4}-\delta\right)
    =\frac{T}{4}-\delta.
\]
Therefore the prior-averaged regret is at least
\begin{align*}
    &\frac{1}{2}
    \left[
        T\left(\frac{1}{4}+\epsilon+\delta\right)
        +\left(\frac{T}{4}-\delta\right)
    \right]
    -T\left(\frac{1}{4}+\epsilon\right)\\
    &\qquad=\frac{T-1}{2}\delta-\frac{T}{2}\epsilon\\
    &\qquad=\frac{w}{2}-\frac{\delta}{8}\\
    &\qquad\ge\frac{3}{8}w.
\end{align*}
By Yao's principle, one of the two deterministic worlds gives at least
this expected regret against the original randomized learner.

For every $s=1,\ldots,T-1$, one has $m_s=c$ and
$|m_T-m_s|=\delta$. Therefore
\[
    \omega_s=\delta.
\]
Hence
\[
    W_T=(T-1)\delta=w.
\]
Finally,
\[
    \mathcal{V}_T=(T-1)\epsilon=\frac{w}{4T}.
\]
This completes the proof.
\end{proof}

\section{The harmonic and primary-price terms are also necessary}
\label{sec:the-harmonic-and-primary-price-terms-are-also-necessary}

The preceding theorem isolates the new secondary-price difficulty. The
two terms inherited from the original problem are independently
unavoidable as well.

\begin{theorem}[Harmonic lower bound]
\label{thm:harmonic-lower-bound}
For every possibly randomized policy and every $T\ge2$, there exists a
deterministic nonincreasing sequence
$p_1\ge\cdots\ge p_T>0$ and the constant secondary-price sequence
$m_t=0$ such that
\[
    \mathcal{V}_T=0,
    \qquad
    W_T=0,
\]
and
\[
    R_T\ge\frac{3}{16}(H_T-1).
\]
\end{theorem}

\emph{Intuition.}
Conditional on the current primary price being $x$, let the next price
remain at $x$ with high probability and fall to
$(1-\epsilon_t)x$ with probability $\epsilon_t$. Uncertainty about
whether the fall occurs costs the learner about $\epsilon_t x$ at round
$t$. The expected erosion of the final-price comparator is only of
order $t\epsilon_t^2x$. Choosing $\epsilon_t=1/(2t)$ makes the net loss
of order $x/t$. Meanwhile, $\sum_t\epsilon_t^2<\infty$, so the expected
price stays bounded away from zero. Summing $1/t$ gives the logarithmic
lower bound.

\begin{proof}
The random construction is only a proof device. Every realization is a
deterministic oblivious sequence.

Set $p_1=1$. Independently for every $t=2,\ldots,T$, conditional on
$p_{t-1}=x$, let
\[
    p_t=
    \begin{cases}
        x & \text{with probability }1-\epsilon_t\\
        (1-\epsilon_t)x & \text{with probability }\epsilon_t
    \end{cases}
    \qquad
    \epsilon_t:=\frac{1}{2t}.
\]
Every realization is nonincreasing. Therefore $M_t=p_t$ and
$\mathcal{V}_T=0$. Since $m_t=0$, we also have $W_T=0$.

Conditional on $p_{t-1}=x$, the expected reward of any offer at round
$t$ is at most
\[
    (1-\epsilon_t)x.
\]
To see this, first condition on a deterministic offer $a$. If
$a<(1-\epsilon_t)x$, it executes in both states and earns at most
$(1-\epsilon_t)x$. If $(1-\epsilon_t)x\le a<x$, it executes only when
there is no drop and its expected reward is
$(1-\epsilon_t)a\le(1-\epsilon_t)x$. If $a\ge x$, it earns zero.
Randomized offers follow by averaging.

Let
\[
    \mu_t:=\mathbb{E}[p_t],
    \qquad
    \mu_0=\mu_1=1,
\]
and set $\epsilon_1=0$. Then
\[
    \mu_t=(1-\epsilon_t^2)\mu_{t-1}.
\]
Moreover,
\[
    \mu_t
    =\prod_{s=2}^t\left(1-\frac{1}{4s^2}\right)
    \ge
    1-\sum_{s=2}^{\infty}\frac{1}{4s^2}
    \ge\frac{3}{4}.
\]

Because the path is nonincreasing, every fixed offer $a<p_T$ executes
on every round. Hence the benchmark satisfies
\[
    B_T\ge Tp_T.
\]
The prior-averaged regret is therefore at least
\[
    T\mu_T-\sum_{t=1}^T(1-\epsilon_t)\mu_{t-1}.
\]
Using
\[
    \sum_{t=1}^T\mu_{t-1}-T\mu_T
    =
    \sum_{t=1}^Tt(\mu_{t-1}-\mu_t)
\]
and $\mu_{t-1}-\mu_t=\epsilon_t^2\mu_{t-1}$, we obtain
\begin{align*}
    \mathbb{E}[R_T]
    &\ge
    \sum_{t=1}^T
    (\epsilon_t-t\epsilon_t^2)\mu_{t-1}\\
    &=
    \frac{1}{4}\sum_{t=2}^T\frac{\mu_{t-1}}{t}\\
    &\ge
    \frac{3}{16}(H_T-1).
\end{align*}
At least one deterministic realization in the support has at least this
regret.
\end{proof}

\begin{theorem}[Primary-price lower bound]
\label{thm:primary-price-lower-bound}
Fix $T\ge2$ and $v\in(0,T-1]$. For every possibly randomized learner,
there exists a deterministic pair
$(\boldsymbol{p},\boldsymbol{m})$ such that
\[
    \mathcal{V}_T\le v,
    \qquad
    W_T=0,
\]
and
\[
    R_T\ge\frac{v}{e}.
\]
\end{theorem}

\begin{proof}
Let
\[
    b:=\frac{v}{T-1}.
\]
Set $m_t=0$ for all $t$, and let
\[
    p_1=\cdots=p_{T-1}=0,
    \qquad
    p_T=X,
\]
where $X\in[0,b]$ has survival function
\[
    \mathbb{P}(X>x)=
    \begin{cases}
        1 & 0\le x<b/e\\
        \dfrac{b}{ex} & b/e\le x<b\\
        0 & x\ge b.
    \end{cases}
\]
Then
\[
    \mathbb{E}[X]=\frac{2b}{e},
\]
and every offer $a$ satisfies
\[
    a\mathbb{P}(X>a)\le\frac{b}{e}.
\]
All learner decisions are made before $X$ is observed, so the learner's
expected reward is at most $Tb/e$. The fixed comparator has payoff
$TX$, whose expectation is $2Tb/e$. Thus the Bayes regret is at least
\[
    \frac{Tb}{e}
    =
    \frac{T}{e(T-1)}v
    \ge
    \frac{v}{e}.
\]
Consequently, one deterministic value of $X$ gives at least this much
regret. For every realization,
\[
    \mathcal{V}_T=(T-1)X\le v,
    \qquad
    W_T=0.
\]
This proves the theorem.
\end{proof}

\end{document}
